% Teil II: durchgerechneter Weg, Teile a--c.

\subsection*{Gegeben und gesucht}
$x\in[\mathrm m]$, $t\in[\mathrm s]$, $v\in[\mathrm{m/s}]$,
$F\in[\mathrm N]$, $m_1,m_2\in[\mathrm{kg}]$; ferner $p$ Druck, $\rho$
Dichte, $g$ Ortsfaktor, $h$ H\"ohe, $x$ Ort, $t$ Zeit. Gesucht: (a) SI-
Einheiten von $C_1,C_2$ in drei Gleichungen, (b) Dimensionsprobe zweier
Formeln, (c) welche Angstroem-Angaben stimmen.

\subsection*{Ansatz}
Die Konstante gleicht die Einheiten aus; Logarithmus-Argument muss
dimensionslos sein; jede Formel beidseitig auf Masse/Laenge/Zeit zurueckfuehren.

\subsection*{Rechnung}
\textbf{a) Einheiten der Konstanten.}
\begin{enumerate}
  \item[(i)] $x=C_1+C_2t$: beide Terme in \unit{\metre}.
\begin{equation}
  C_1:\ \unit{\metre},\qquad
  C_2:\ \frac{\unit{\metre}}{\unit{\second}}.
\end{equation}
  \item[(ii)] $v=C_1\,x\,\ln(C_2t)$: $\ln$-Argument dimensionslos
        $\Rightarrow C_2$ in \unit{\per\second}.
\begin{equation}
  C_1=\frac{v}{x}=\frac{\unit{\metre\per\second}}{\unit{\metre}}
      =\unit{\per\second},\qquad
  C_2:\ \unit{\per\second}.
\end{equation}
  \item[(iii)] $F=C_1\,\frac{m_1m_2}{x^2}$:
\begin{equation}
  C_1=\frac{F\,x^2}{m_1m_2}
      =\frac{\unit{\newton\metre\squared}}{\unit{\kilogram\squared}}
      =\frac{\unit{\metre\cubed}}{\unit{\kilogram\second\squared}}
      \quad(\text{Gravitationskonstante }G)\;.
\end{equation}
\end{enumerate}

\textbf{b) Dimensionsprobe.}
\begin{enumerate}
  \item[(i)] $p+\rho gh+\tfrac12\rho v^2=\text{const.}$ Alle drei Terme:
\begin{equation}
  \frac{\unit{\kilogram}}{\unit{\metre\second\squared}}
  \quad\text{(Druck, Bernoulli)}\;.
\end{equation}
$p=\unit{\pascal}$: ok. $\rho gh$: $\tfrac{\mathrm{kg}}{\mathrm m^3}\cdot
\frac{\mathrm m}{\mathrm s^2}\cdot\mathrm m=\frac{\mathrm{kg}}{\mathrm{m s^2}}$. ok.
$\tfrac12\rho v^2$: $\frac{\mathrm{kg}}{\mathrm m^3}\cdot
\frac{\mathrm m^2}{\mathrm s^2}=\frac{\mathrm{kg}}{\mathrm{m s^2}}$. ok.
Die Formel kann stimmen.

  \item[(ii)] $F=\mathrm d^2x/\mathrm dt^2$: links
        \unit{\newton}$=\frac{\mathrm{kg\,m}}{\mathrm s^2}$, rechts
\begin{equation}
  \frac{\mathrm d^2x}{\mathrm dt^2}:\ \frac{\unit{\metre}}{\unit{\second\squared}},
\end{equation}
also fehlt die Masse. Die Formel \textbf{stimmt nicht} (richtig:
$F=m\,\mathrm d^2x/\mathrm dt^2$).
\end{enumerate}

\textbf{c) Angstroem-Umrechnung.} $1\,\mathrm{\AA}=10^{-10}\mathrm m$:
\begin{align}
  (i)\ 10\,\mathrm{nm}&=10^{-8}\mathrm m\ \text{(falsch)},\\
  (ii)\ 10^5\,\mathrm{fm}&=10^5\cdot10^{-15}\mathrm m=10^{-10}\mathrm m\
      \checkmark,\\
  (iii)\ 10^{-1}\,\mathrm{nm}&=10^{-1}\cdot10^{-9}\mathrm m=10^{-10}\mathrm m\
      \checkmark,\\
  (iv)\ 10^{-5}\,\mathrm{fm}&=10^{-20}\mathrm m\ \text{(falsch)},\\
  (v)\ 10^{-4}\,\mathrm{\mu m}&=10^{-4}\cdot10^{-6}\mathrm m=10^{-10}\mathrm m\
      \checkmark,\\
  (vi)\ 10^{3}\,\mathrm{\mu m}&=10^{-3}\mathrm m\ \text{(falsch)}.
\end{align}

\begin{ergebnis}
(a) $C_1$: m (bzw. s$^{-1}$, bzw. m$^3$/(kg$\,\cdot$s$^2$)); $C_2$: m/s
(bzw. s$^{-1}$). (b) Bernoulli-Gleichung ok, $F=\mathrm d^2x/\mathrm dt^2$
nicht. (c) Richtig: (ii), (iii), (v).
\end{ergebnis}

\begin{kontrolle}
\textbf{Einheiten:} in (a) jede Gleichung links/rechts auf eine
Einheitenkette geprueft; in (c) alle auf \unit{\metre}. \checkmark

\textbf{Querprobe:} Durchmesser eines Atoms $\sim1\mathrm{\AA}$.
$10^5\,\mathrm{fm}=10^{-10}\mathrm m$ ist der gleiche Wert, also sind
(ii),(iii),(v) konsistent. \checkmark
\end{kontrolle}

\begin{merksatz}
Eine Konstante in einer Formel uebernimmt genau die Einheit, die die
Gleichung dimensionell ausgleicht; naehert man eine Gleichung auf die
Dimension aller Terme, erkennt man fehlende Faktoren (hier die Masse) und
kann Umrechnungen mit 10er-Potenzen sauber pruefen.
\end{merksatz}
